% !TeX root = ../papers/nc-p3.tex
\section{反対行列の一列の公因子と七余因子表示での接触}
\label{sec:opposite-column-divisor}
\begin{theorem}\label{thm:opposite-column-divisor}
四辺有限の標準有限Frobenius窓で、全二次式を成分に持つ反対平方根 $N$ は一般点で可逆とする。
任意の非零定数列 $v$ について、$N(x)v$ の成分に公因子があれば
\[
                    N(x)v=l(x) Bx,\qquad v^t Bx\in k^\times l(x)
\]
となる可逆行列 $B$ と一次式 $l$ がある。この場合 $l^2$ は行列式を割る。
特に点窓が被約ならどの列にも公因子はなく、一般核の入射の全ての出力繊維は次元1以下、
その曲線部分の次数は重複を込めて2以下である。階数2軌跡の次元や八次式の整性を仮定しない。
\end{theorem}
\begin{proof}[残余二次評価の単元作用]
$N(x)v$ は独立な四つの二次式なので公因子は一次式で商の四一次式は独立である。
入力を $Bx$ と変更し $B=I$ とする。$m=v^tx$ と置く。点窓の有限環で諸線束を自明化し
$N=ww^t$ と書く（単元倍率の平方根を取る）。$x,w$ はともに単位イデアルを生成する列で
$(v^tw)w=lx$ だから $(v^tw)=(l)$。局所Artin環の積で二生成元は単元倍となるので $v^tw=b l$ とできる。
$D=H^0(X,\mathcal O_X)/\operatorname{ann}(l)$ とすればこの残余上で $bw=x$。
$D$ には $\beta=b^{-2}$ により
\[
                    N(x)=\beta xx^t,\qquad l=\beta m
\]
が成立する。$D$ の長さは $20-\ell(X\cap V(l))\ge13$ であり、二次式の制限は単射である。
実際制限が零なら元の窓で $l$ 倍が零となり三次評価の単射性を使える。
従って二次評価 $Q\subset D$ は十次元で、十個の行列成分の恒等式から単元 $\beta$ による乗法で保たれる。

$l,m$ が独立と仮定する。$\beta$ の異なる固有値 $\lambda$ ごとに $D_\lambda$ へ分け、
$t=\beta-\lambda$ による冪濾過の非零層 $t^j D_\lambda/t^{j+1}D_\lambda$ を取る。
各層は $1\mapsto t^j$ により $D_\lambda/(\operatorname{ann}(t^j)+(t))$ と同一視できる。
従って $l-\lambda m=0$ の平面上の、長さ $a\le7$ の部分スキームの関数空間である。
各層への二次式の評価階数を $h_2$ とする。$Q$ は固有射影と $t$ 倍で保たれるので、
$Q$ の誘導濾過の当該次数商の次元は $h_2$ 以上である（射影した二次評価の $t^j$ 倍を使う）。
全層の $h_2$ の和は10以下となる。

長さ $a\le5$ なら $h_2=a$ である。長さ2以上なら長さ2以上を切る直線を選ぶ
（長さ2の部分スキームは共線）。切断長3が可能ならそれを選び、残余長は2以下で一次式に独立な条件を課す。
最大切断長2なら残余は長さ3以下で、長さ3の場合は非共線である。いずれも直線の残余完全列で
二次式が全射となる。直線切断長は常に3以下である。$a\ge6$ のとき $h_2\ge5$ となる。
実際独立な二平面二次式が共通因子を持たなければ長さ4で尽き、共通直線があればその残余は
二本の独立な一次式で消され長さ1以下だからである。
層が三つ以上なら、長さ6以上の層が二つ以上のとき階数和は少なくとも $5+5+1=11$。
そのような層が高々一つでも全長からの損失は高々2で、やはり $13-2>10$。
従って層は高々二つである。この層数は $\beta$ の最小多項式の次数であるから
$\prod_\lambda (l-\lambda m)^{\nu_\lambda}$ は $D$ 上で消える次数2以下の非零多項式になる
（$\nu_\lambda$ は各冪零指数）。二次評価の単射性に反する。

従って $m=c l$。入力と出力の対を合わせた基底で $v$ を第一座標、$l=x_0$ に規格化すれば、
定数倍を除いて第一列は $x_0 x$ に等しい。第一行と第一列から $x_0$ をくくれる
（左上は $x_0^2$）ので $l^2\mid\det N$ である。
被約なら $D$ は $l\ne0$ の印の直積で $\beta$ は一定値となる。二次評価の単射性により
$N$ は多項式として $xx^t$ の定数倍になり、一般可逆性に反する。
\end{proof}
\begin{proof}[固定核の曲線次数]
被約の場合は四つの列成分の一般二式が二次二次完全交叉曲線 $Y$ を切る。
共通零スキームの純一次元部分 $C$（孤立・埋込零次元部分を除く）の次数が4なら $C=Y$ であり
四つの独立な二次式を持てない。次数3なら $Y$ 内の残余はCohen--Macaulayな次数1の直線 $J$。
完全交叉リンケージの $\mathcal I_{C/Y}=\omega_J\otimes\omega_Y^{-1}$ と $\omega_Y=\mathcal O_Y$ より
$h^0(\mathcal I_C(2))\le2+h^0(\omega_J(2))=3$ で、これも不可能である。
\end{proof}

\begin{corollary}\label{cor:bezout-contact-divisor}
$N(x)=\operatorname{Bez}(T(sx),T(tx))$ が可逆な $T:S_4\to S_4$ から来る上記の反対根とする。
被約窓では二つの有理正規四次曲線 $R,T^*R$ は交わらない。
非被約でも交点では必ず両接線が一致する。この結論に結果式の整性は要らない。
\end{corollary}
\begin{proof}
$T^*\mathrm{ev}_z$ が $\mathrm{ev}_w$ に比例するとし、$w,z$ で消える一次式を $a_w,a_z$ とする。
Bezout双一次式の片方に $z$ を入れると、二変数の交代分子を対角の式で割る公式から、
得る列は定数倍を除いて
\[
                   x(w)\,B_z x,\qquad B_z x=T(a_w x)/a_z
\]
である。$T$ は $a_w S_3$ を $a_z S_3$ に同型に運ぶので $B_z$ は可逆である。
前定理が被約の場合を除く。一般の場合でも $(B_z x)(z)$ は $x(w)$ に比例する。
左辺は $T(a_w x)$ の $z$ での一階微分の定数倍であり、この条件は
$T^*\mathrm{ev}'_z\in\langle\mathrm{ev}_w,\mathrm{ev}'_w\rangle$ と同値である。
非退化な埋込みの接方向なので両接線は一致する。
\end{proof}
